\documentclass[12pt, a4paper]{ctexart}
\usepackage{fontspec}
\usepackage{xcolor}
\usepackage{hyperref}
\usepackage{geometry}
\usepackage{enumitem}
\usepackage{parskip}
\usepackage{titlesec}
\usepackage{tcolorbox}
\tcbuselibrary{breakable}
\usepackage{setspace}
\usepackage{listings}

\geometry{left=2.5cm, right=2.5cm, top=2.5cm, bottom=2.5cm}

\setmainfont{Times New Roman}
\setCJKmainfont{SimSun}

\hypersetup{
    colorlinks=true,
    linkcolor=blue!70!black,
    urlcolor=cyan,
    pdftitle={Java API-逐题精讲解义},
    pdfauthor={专升本-fifine老师}
}

\definecolor{myblue}{RGB}{30,100,200}
\definecolor{lightblue}{RGB}{230,240,255}
\definecolor{darkblue}{RGB}{0,50,120}

\newtcolorbox{bluebox}[1][]{
    colback=lightblue,
    colframe=myblue,
    coltitle=darkblue,
    fonttitle=\bfseries,
    title={#1},
    boxrule=0.8pt,
    arc=4pt,
    left=6pt,
    right=6pt,
    top=6pt,
    bottom=6pt,
    before skip=8pt,
    after skip=8pt,
    breakable
}

\usepackage{etoolbox}
\BeforeBeginEnvironment{bluebox}{\par\vfill}%
\AfterEndEnvironment{bluebox}{\par\vfill}%

\titleformat{\section}
  {\normalfont\Large\bfseries\color{myblue}}{\thesection}{1em}{}
\titlespacing*{\section}{0pt}{3.5ex plus 1ex minus .2ex}{1.5ex plus .2ex}

\title{\textcolor{myblue}{\textbf{Java API - 逐题精讲解义}}}
\author{专升本-fifine老师 教学组}
\date{2026年03月02日}

\lstset{
    language=Java,
    basicstyle=\ttfamily\small,
    keywordstyle=\color{blue},
    commentstyle=\color{green!60!black},
    stringstyle=\color{orange},
    numbers=left,
    numberstyle=\tiny\color{gray},
    frame=single,
    breaklines=true,
    columns=fullflexible,
    showstringspaces=false
}

\newcommand{\code}[1]{\texttt{#1}}

\begin{document}

\maketitle
\thispagestyle{empty}
\newpage

\section{Java API}

\begin{enumerate}[label=\textbf{\arabic*.}]

	\item \textbf{以下哪个方法可以将字符串转换为整数？}
	      \begin{enumerate}[label=\Alph*.]
		      \item parseInt()
		      \item valueOf()
		      \item toString()
		      \item toInteger()
	      \end{enumerate}

	      \begin{bluebox}{答案与深度解析}
		      \textbf{答案：A. parseInt()}

		      \textbf{【核心认知框架】}

		      \textbf{【选项原理深度拆解】}
					  \begin{itemize}[label={}, leftmargin=*, nosep]


		      \item \textbf{A. parseInt() — 正确}
					  \begin{itemize}[label={}, nosep]
		      
			      \item Integer.parseInt(String s) 将字符串解析为int基本类型
			      \item 属于Integer类的静态方法

					  \end{itemize}
		      \item \textbf{B. valueOf() — 错误}
					  \begin{itemize}[label={}, nosep]
		      
			      \item Integer.valueOf(String s) 返回Integer包装类对象
			      \item 返回类型是Integer而非int
		      

					  \end{itemize}
		      \item \textbf{C. toString() — 错误}
					  \begin{itemize}[label={}, nosep]
		      
			      \item toString()是将其他类型转为字符串，不是字符串转整数
		      

					  \end{itemize}
		      \item \textbf{D. toInteger() — 错误}
					  \begin{itemize}[label={}, nosep]
		      
			      \item 不存在此方法
		      

		      
					  \end{itemize}
					  \end{itemize}
\textbf{【应背诵知识点】}

		      parseInt返回int基本类型，valueOf返回Integer包装类
		      \end{bluebox}

	      \vspace{1em}

	\item \textbf{在Java中，以下哪个方法用于获取字符串的长度？}
	      \begin{enumerate}[label=\Alph*.]
		      \item length()
		      \item size()
		      \item count()
		      \item getLength()
	      \end{enumerate}

	      \begin{bluebox}{答案与深度解析}
		      \textbf{答案：A. length()}

		      \textbf{【选项原理深度拆解】}
					  \begin{itemize}[label={}, leftmargin=*, nosep]


		      \item \textbf{A. length() — 正确}
					  \begin{itemize}[label={}, nosep]
		      
			      \item String类使用length()方法获取字符数
		      

					  \end{itemize}
		      \item \textbf{B. size() — 错误}
					  \begin{itemize}[label={}, nosep]
		      
			      \item size()用于集合类（如List、Set）

					  \end{itemize}
		      \item \textbf{C. count() — 错误}
					  \begin{itemize}[label={}, nosep]
		      
			      \item 不是String的方法
		      

					  \end{itemize}
		      \item \textbf{D. getLength() — 错误}
					  \begin{itemize}[label={}, nosep]
		      
			      \item 不是String的标准方法
		      
		      

	      
					  \end{itemize}
					  \end{itemize}
\end{bluebox}

	      \vspace{1em}

	\item \textbf{以下哪个方法用于将字符串转换为大写？}
	      \begin{enumerate}[label=\Alph*.]
		      \item toUpperCase()
		      \item upperCase()
		      \item makeUpperCase()
		      \item convertToUpperCase()
	      \end{enumerate}

	      \begin{bluebox}{答案与深度解析}
		      \textbf{答案：A. toUpperCase()}

		      \textbf{【选项原理深度拆解】}
					  \begin{itemize}[label={}, leftmargin=*, nosep]


		      \item \textbf{A. toUpperCase() — 正确}
					  \begin{itemize}[label={}, nosep]
		      
			      \item String类的toUpperCase()方法返回大写字符串

					  \end{itemize}
		      \item \textbf{B/C/D — 错误}
					  \begin{itemize}[label={}, nosep]
		      
			      \item 这些方法不存在
		      
		      

	      
					  \end{itemize}
					  \end{itemize}
\end{bluebox}

	      \vspace{1em}

	\item \textbf{以下哪个方法用于将字符串转换为小写？}
	      \begin{enumerate}[label=\Alph*.]
		      \item toLowerCase()
		      \item lowerCase()
		      \item makeLowerCase()
		      \item convertToLowerCase()
	      \end{enumerate}

	      \begin{bluebox}{答案与深度解析}
		      \textbf{答案：A. toLowerCase()}

		      \textbf{【选项原理深度拆解】}
					  \begin{itemize}[label={}, leftmargin=*, nosep]


		      \item \textbf{A. toLowerCase() — 正确}
					  \begin{itemize}[label={}, nosep]
		      
			      \item String类的toLowerCase()方法返回小写字符串

					  \end{itemize}
		      \item \textbf{B/C/D — 错误}
					  \begin{itemize}[label={}, nosep]
		      
			      \item 这些方法不存在
		      
		      

	      
					  \end{itemize}
					  \end{itemize}
\end{bluebox}

	      \vspace{1em}

	\item \textbf{以下哪个方法用于截取字符串？}
	      \begin{enumerate}[label=\Alph*.]
		      \item substring()
		      \item cut()
		      \item slice()
		      \item truncate()
	      \end{enumerate}

	      \begin{bluebox}{答案与深度解析}
		      \textbf{答案：A. substring()}

		      \textbf{【选项原理深度拆解】}
					  \begin{itemize}[label={}, leftmargin=*, nosep]


		      \item \textbf{A. substring() — 正确}
					  \begin{itemize}[label={}, nosep]
		      
			      \item substring(int beginIndex, int endIndex) 截取字符串
			      \item 注意：endIndex是Exclusive的

					  \end{itemize}
		      \item \textbf{B/C/D — 错误}
					  \begin{itemize}[label={}, nosep]
		      
			      \item 这些方法不是String类的
		      
		      

	      
					  \end{itemize}
					  \end{itemize}
\end{bluebox}

	      \vspace{1em}

	\item \textbf{以下哪个方法用于判断字符串是否以指定字符串开头？}
	      \begin{enumerate}[label=\Alph*.]
		      \item startsWith()
		      \item beginWith()
		      \item isStartWith()
		      \item hasStartWith()
	      \end{enumerate}

	      \begin{bluebox}{答案与深度解析}
		      \textbf{答案：A. startsWith()}

		      \textbf{【选项原理深度拆解】}
					  \begin{itemize}[label={}, leftmargin=*, nosep]


		      \item \textbf{A. startsWith() — 正确}
					  \begin{itemize}[label={}, nosep]
		      
			      \item boolean startsWith(String prefix) 判断是否以指定字符串开头

					  \end{itemize}
		      \item \textbf{B/C/D — 错误}
					  \begin{itemize}[label={}, nosep]
		      
			      \item 这些方法不存在
		      
		      

	      
					  \end{itemize}
					  \end{itemize}
\end{bluebox}

	      \vspace{1em}

	\item \textbf{以下哪个方法用于判断字符串是否以指定字符串结尾？}
	      \begin{enumerate}[label=\Alph*.]
		      \item endsWith()
		      \item finishWith()
		      \item isEndWith()
		      \item hasEndWith()
	      \end{enumerate}

	      \begin{bluebox}{答案与深度解析}
		      \textbf{答案：A. endsWith()}

		      \textbf{【选项原理深度拆解】}
					  \begin{itemize}[label={}, leftmargin=*, nosep]


		      \item \textbf{A. endsWith() — 正确}
					  \begin{itemize}[label={}, nosep]
		      
			      \item boolean endsWith(String suffix) 判断是否以指定字符串结尾

					  \end{itemize}
		      \item \textbf{B/C/D — 错误}
					  \begin{itemize}[label={}, nosep]
		      
			      \item 这些方法不存在
		      
		      
			  
	      
					  \end{itemize}
					  \end{itemize}
\end{bluebox}

	      \vspace{1em}

	\item \textbf{以下哪个方法用于在字符串中查找指定字符第一次出现的位置？}
	      \begin{enumerate}[label=\Alph*.]
		      \item indexOf()
		      \item find()
		      \item search()
		      \item locate()
	      \end{enumerate}

	      \begin{bluebox}{答案与深度解析}
		      \textbf{答案：A. indexOf()}

		      \textbf{【选项原理深度拆解】}
					  \begin{itemize}[label={}, leftmargin=*, nosep]


		      \item \textbf{A. indexOf() — 正确}
					  \begin{itemize}[label={}, nosep]
		      
			      \item int indexOf(String str) 返回首次出现的索引
			      \item 未找到返回-1

					  \end{itemize}
		      \item \textbf{B/C/D — 错误}
					  \begin{itemize}[label={}, nosep]
		      
			      \item 这些方法不存在于String类
		      
		      

	      
					  \end{itemize}
					  \end{itemize}
\end{bluebox}

	      \vspace{1em}

	\item \textbf{以下哪个方法用于在字符串中查找指定字符最后一次出现的位置？}
	      \begin{enumerate}[label=\Alph*.]
		      \item lastIndexOf()
		      \item findLast()
		      \item searchLast()
		      \item locateLast()
	      \end{enumerate}

	      \begin{bluebox}{答案与深度解析}
		      \textbf{答案：A. lastIndexOf()}

		      \textbf{【选项原理深度拆解】}
					  \begin{itemize}[label={}, leftmargin=*, nosep]


		      \item \textbf{A. lastIndexOf() — 正确}
					  \begin{itemize}[label={}, nosep]
		      
			      \item int lastIndexOf(String str) 返回最后一次出现的索引
			      \item 未找到返回-1

					  \end{itemize}
		      \item \textbf{B/C/D — 错误}
					  \begin{itemize}[label={}, nosep]
		      
			      \item 这些方法不存在于String类
		      
		      

	      
					  \end{itemize}
					  \end{itemize}
\end{bluebox}

	      \vspace{1em}

	\item \textbf{以下关于Java中的包装类的说法错误的是}
	      \begin{enumerate}[label=\Alph*.]
		      \item 包装类可以将基本数据类型转换为对象
		      \item Integer是int的包装类
		      \item 包装类的对象不可变
		      \item 包装类可以提高基本数据类型的操作效率
	      \end{enumerate}

	      \begin{bluebox}{答案与深度解析}
		      \textbf{答案：D. 包装类可以提高基本数据类型的操作效率}

		      \textbf{【核心认知框架】}

		      \textbf{【选项原理深度拆解】}
					  \begin{itemize}[label={}, leftmargin=*, nosep]


		      \item \textbf{A. 包装类可以将基本数据类型转换为对象 — 正确}
					  \begin{itemize}[label={}, nosep]
		      
			      \item 装箱(Autoboxing)：int → Integer

					  \end{itemize}
		      \item \textbf{B. Integer是int的包装类 — 正确}
					  \begin{itemize}[label={}, nosep]
		      
			      \item 八大基本类型都有对应包装类
		      

					  \end{itemize}
		      \item \textbf{C. 包装类的对象不可变 — 正确}
					  \begin{itemize}[label={}, nosep]
		      
			      \item Integer等包装类都是不可变类（immutable）
		      

					  \end{itemize}
		      \item \textbf{D. 包装类可以提高基本数据类型的操作效率 — 错误}
					  \begin{itemize}[label={}, nosep]
		      
			      \item 恰恰相反，包装类有额外开销，效率更低
			      \item 包装类主要用于集合操作、泛型等需要对象的场景
		      

		      
					  \end{itemize}
					  \end{itemize}
\textbf{【应背诵知识点】}

		      包装类特点：
		      \begin{itemize}[label={}, nosep]
			      \item 基本类型与包装类的转换（装箱/拆箱）
			      \item 自动装箱：Integer i = 10;（JDK5+）
			      \item 不可变性保证线程安全
			      \item 缓存机制：Integer IntegerCache [-128, 127]
		      \end{itemize}
		      \end{bluebox}

	      \vspace{1em}

	\item \textbf{以下哪个是Double类的静态方法，用于将字符串转换为Double类型？}
	      \begin{enumerate}[label=\Alph*.]
		      \item parseDouble()
		      \item valueOf()
		      \item toString()
		      \item toDouble()
	      \end{enumerate}

	      \begin{bluebox}{答案与深度解析}
		      \textbf{答案：A. parseDouble()}

		      \textbf{【选项原理深度拆解】}
					  \begin{itemize}[label={}, leftmargin=*, nosep]


		      \item \textbf{A. parseDouble() — 正确}
					  \begin{itemize}[label={}, nosep]
		      
			      \item Double.parseDouble(String s) 返回double基本类型

					  \end{itemize}
		      \item \textbf{B. valueOf() — 错误}
					  \begin{itemize}[label={}, nosep]
		      
			      \item Double.valueOf(String s) 返回Double对象
		      

					  \end{itemize}
		      \item \textbf{C/D — 错误}
					  \begin{itemize}[label={}, nosep]
		      
			      \item 这些方法不是Double类的转换方法
		      
		      

	      
					  \end{itemize}
					  \end{itemize}
\end{bluebox}

	      \vspace{1em}

	\item \textbf{以下关于Java中的枚举类型的说法错误的是}
	      \begin{enumerate}[label=\Alph*.]
		      \item 枚举类型是一种特殊的类
		      \item 枚举类型的成员是常量
		      \item 枚举类型可以有构造方法
		      \item 枚举类型不能实现接口
	      \end{enumerate}

	      \begin{bluebox}{答案与深度解析}
		      \textbf{答案：D. 枚举类型不能实现接口}

		      \textbf{【核心认知框架】}

		      \textbf{【选项原理深度拆解】}
					  \begin{itemize}[label={}, leftmargin=*, nosep]


		      \item \textbf{A. 枚举类型是一种特殊的类 — 正确}
					  \begin{itemize}[label={}, nosep]
		      
			      \item 枚举本质上继承自java.lang.Enum

					  \end{itemize}
		      \item \textbf{B. 枚举类型的成员是常量 — 正确}
					  \begin{itemize}[label={}, nosep]
		      
			      \item 枚举成员是public static final的实例
		      

					  \end{itemize}
		      \item \textbf{C. 枚举类型可以有构造方法 — 正确}
					  \begin{itemize}[label={}, nosep]
		      
			      \item 枚举构造方法必须是private
		      

					  \end{itemize}
		      \item \textbf{D. 枚举类型不能实现接口 — 错误}
					  \begin{itemize}[label={}, nosep]
		      
			      \item 枚举可以实现接口！
			      \item 格式：enum Color implements Runnable \{ RED \{ @Override public void run() \{\} \} \}
		      

		      
					  \end{itemize}
					  \end{itemize}
\textbf{【应背诵知识点】}

		      枚举特点：
		      \begin{itemize}[label={}, nosep]
			      \item 继承java.lang.Enum
			      \item 成员是static final实例
			      \item 构造方法私有化
			      \item 可以实现接口
			      \item 可以有抽象方法（每个实例必须实现）
		      \end{itemize}
		      \end{bluebox}

	      \vspace{1em}

	\item \textbf{以下哪个类用于获取系统属性？}
	      \begin{enumerate}[label=\Alph*.]
		      \item System
		      \item Properties
		      \item Runtime
		      \item SystemProperties
	      \end{enumerate}

	      \begin{bluebox}{答案与深度解析}
		      \textbf{答案：B. Properties}

		      \textbf{【选项原理深度拆解】}
					  \begin{itemize}[label={}, leftmargin=*, nosep]


		      \item \textbf{A. System — 错误}
					  \begin{itemize}[label={}, nosep]
		      
			      \item System类是系统类，通过getProperties()获取属性

					  \end{itemize}
		      \item \textbf{B. Properties — 正确}
					  \begin{itemize}[label={}, nosep]
		      
			      \item System.getProperties()返回Properties对象
			      \item Properties继承自Hashtable
		      

					  \end{itemize}
		      \item \textbf{C. Runtime — 错误}
					  \begin{itemize}[label={}, nosep]
		      
			      \item Runtime是运行时类，用于内存管理等
		      

					  \end{itemize}
		      \item \textbf{D. SystemProperties — 错误}
					  \begin{itemize}[label={}, nosep]
		      
			      \item 没有这个类
		      
		      

	      
					  \end{itemize}
					  \end{itemize}
\end{bluebox}

	      \vspace{1em}

	\item \textbf{以下关于Java中的注解的说法错误的是}
	      \begin{enumerate}[label=\Alph*.]
		      \item 注解是一种元数据
		      \item 可以自定义注解
		      \item 注解可以影响程序的逻辑
		      \item 注解可以为代码添加额外的信息
	      \end{enumerate}

	      \begin{bluebox}{答案与深度解析}
		      \textbf{答案：C. 注解可以影响程序的逻辑}

		      \textbf{【核心认知框架】}

		      \textbf{【选项原理深度拆解】}
					  \begin{itemize}[label={}, leftmargin=*, nosep]


		      \item \textbf{A. 注解是一种元数据 — 正确}
					  \begin{itemize}[label={}, nosep]
		      
			      \item 注解用于提供程序的元数据信息

					  \end{itemize}
		      \item \textbf{B. 可以自定义注解 — 正确}
					  \begin{itemize}[label={}, nosep]
		      
			      \item 使用@interface定义
		      

					  \end{itemize}
		      \item \textbf{C. 注解可以影响程序的逻辑 — 错误}
					  \begin{itemize}[label={}, nosep]
		      
			      \item 注解本身不改变程序逻辑
			      \item 需要由注解处理器（工具/框架）处理
		      

					  \end{itemize}
		      \item \textbf{D. 注解可以为代码添加额外的信息 — 正确}
					  \begin{itemize}[label={}, nosep]
		      
			      \item 这是注解的主要作用
		      
		      

	      
					  \end{itemize}
					  \end{itemize}
\end{bluebox}

	      \vspace{1em}

	\item \textbf{以下哪个是Java中的基本注解？}
	      \begin{enumerate}[label=\Alph*.]
		      \item @Override
		      \item @Deprecated
		      \item @SuppressWarnings
		      \item 以上都是
	      \end{enumerate}

	      \begin{bluebox}{答案与深度解析}
		      \textbf{答案：D. 以上都是}

		      \textbf{【核心认知框架】}

		      \textbf{【选项原理深度拆解】}
					  \begin{itemize}[label={}, leftmargin=*, nosep]


		      \item[] 三大内置注解：
		      
			      \item @Override：标记重写方法
			      \item @Deprecated：标记已过时
			      \item @SuppressWarnings：抑制警告
		      
	      
					  \end{itemize}
\end{bluebox}

	      \vspace{1em}

	\item \textbf{以下关于Java注解@SuppressWarnings的作用，正确的是}
	      \begin{enumerate}[label=\Alph*.]
		      \item 抑制所有的警告
		      \item 抑制指定类型的警告
		      \item 增加警告信息
		      \item 用于代码注释
	      \end{enumerate}

	      \begin{bluebox}{答案与深度解析}
		      \textbf{答案：B. 抑制指定类型的警告}

		      \textbf{【选项原理深度拆解】}
					  \begin{itemize}[label={}, leftmargin=*, nosep]


		      \item \textbf{A. 抑制所有的警告 — 错误}
					  \begin{itemize}[label={}, nosep]
		      
			      \item 需要指定具体警告类型

					  \end{itemize}
		      \item \textbf{B. 抑制指定类型的警告 — 正确}
					  \begin{itemize}[label={}, nosep]
		      
			      \item @SuppressWarnings("unchecked")、@SuppressWarnings("deprecation")等
		      

					  \end{itemize}
		      \item \textbf{C/D — 错误}
					  \begin{itemize}[label={}, nosep]
		      
			      \item 与注释无关，不能增加警告
		      
		      

	      
					  \end{itemize}
					  \end{itemize}
\end{bluebox}

	      \vspace{1em}

	\item \textbf{以下关于Java中StringBuilder和StringBuffer的说法正确的是}
	      \begin{enumerate}[label=\Alph*.]
		      \item StringBuilder是线程安全的，StringBuffer不是
		      \item StringBuffer是线程安全的，StringBuilder不是
		      \item 两者性能相同
		      \item 两者都不是线程安全的
	      \end{enumerate}

	      \begin{bluebox}{答案与深度解析}
		      \textbf{答案：B. StringBuffer是线程安全的，StringBuilder不是}

		      \textbf{【核心认知框架】}

		      \textbf{【选项原理深度拆解】}
					  \begin{itemize}[label={}, leftmargin=*, nosep]


		      \item \textbf{A. StringBuilder线程安全 — 错误}
					  \begin{itemize}[label={}, nosep]
		      
			      \item StringBuilder是非线程安全的
		      

					  \end{itemize}
		      \item \textbf{B. StringBuffer线程安全 — 正确}
					  \begin{itemize}[label={}, nosep]
		      
			      \item StringBuffer方法使用synchronized修饰
			      \item 线程安全，但性能较低
		      

					  \end{itemize}
		      \item \textbf{C. 性能相同 — 错误}
					  \begin{itemize}[label={}, nosep]
		      
			      \item StringBuilder性能更好（无需同步）
		      

					  \end{itemize}
		      \item \textbf{D. 都不是线程安全 — 错误}
					  \begin{itemize}[label={}, nosep]
		      
			      \item StringBuffer是线程安全的
		      

		      
					  \end{itemize}
					  \end{itemize}
\textbf{【应背诵知识点】}

		      StringBuilder vs StringBuffer：
		      \begin{itemize}[label={}, nosep]
			      \item StringBuilder：非线程安全，性能好
			      \item StringBuffer：线程安全（JDK1.0），性能较差
			      \item 单线程用StringBuilder，多线程用StringBuffer
			      \item String：不可变，线程安全
		      \end{itemize}
	      \end{bluebox}

	      \vspace{1em}

	\item \textbf{以下哪个方法用于将StringBuilder或StringBuffer对象转换为字符串？}
	      \begin{enumerate}[label=\Alph*.]
		      \item toString()
		      \item toStr()
		      \item getString()
		      \item getText()
	      \end{enumerate}

	      \begin{bluebox}{答案与深度解析}
		      \textbf{答案：A. toString()}

		      \textbf{【选项原理深度拆解】}
					  \begin{itemize}[label={}, leftmargin=*, nosep]


		      \item \textbf{A. toString() — 正确}
					  \begin{itemize}[label={}, nosep]
		      
			      \item Object类的方法，所有对象都有
			      \item StringBuilder/StringBuffer都重写了toString()
		      

					  \end{itemize}
		      \item \textbf{B/C/D — 错误}
					  \begin{itemize}[label={}, nosep]
		      
			      \item 这些方法不存在
		      
	      
					  \end{itemize}
					  \end{itemize}
\end{bluebox}

	      \vspace{1em}

	\item \textbf{以下关于Java中Math类的说法错误的是}
	      \begin{enumerate}[label=\Alph*.]
		      \item 提供了常用的数学计算方法
		      \item 所有方法都是静态的
		      \item 可以创建Math类的对象
		      \item 包含了求绝对值、平方根等方法
	      \end{enumerate}

	      \begin{bluebox}{答案与深度解析}
		      \textbf{答案：C. 可以创建Math类的对象}

		      \textbf{【核心认知框架】}

		      \textbf{【选项原理深度拆解】}
					  \begin{itemize}[label={}, leftmargin=*, nosep]


		      \item \textbf{A. 提供数学计算方法 — 正确}
					  \begin{itemize}[label={}, nosep]
		      
			      \item Math类提供丰富的数学函数
		      

					  \end{itemize}
		      \item \textbf{B. 所有方法都是静态的 — 正确}
					  \begin{itemize}[label={}, nosep]
		      
			      \item Math是工具类，方法都是静态的
		      

					  \end{itemize}
		      \item \textbf{C. 可以创建Math类的对象 — 错误}
					  \begin{itemize}[label={}, nosep]
		      
			      \item Math构造方法是private，不能实例化
		      

					  \end{itemize}
		      \item \textbf{D. 包含绝对值、平方根 — 正确}
					  \begin{itemize}[label={}, nosep]
		      
			      \item Math.abs()、Math.sqrt()等
		      

		      
					  \end{itemize}
					  \end{itemize}
\textbf{【应背诵知识点】}

		      Math类特点：
		      \begin{itemize}[label={}, nosep]
			      \item 工具类：构造方法私有
			      \item 所有方法static：直接Math.method()
			      \item 常量：Math.PI、Math.E
			      \item 常见方法：abs、sqrt、pow、random、max、min
		      \end{itemize}
	      \end{bluebox}

	      \vspace{1em}

	\item \textbf{以下哪个方法用于获取随机数？}
	      \begin{enumerate}[label=\Alph*.]
		      \item random()
		      \item getRandom()
		      \item Random.nextInt()
		      \item Math.random()
	      \end{enumerate}

	      \begin{bluebox}{答案与深度解析}
		      \textbf{答案：D. Math.random()}

		      \textbf{【选项原理深度拆解】}
					  \begin{itemize}[label={}, leftmargin=*, nosep]


		      \item \textbf{A. random() — 错误}
					  \begin{itemize}[label={}, nosep]
		      
			      \item 不是任何类的直接方法
		      

					  \end{itemize}
		      \item \textbf{B. getRandom() — 错误}
					  \begin{itemize}[label={}, nosep]
		      
			      \item 不存在此方法
		      

					  \end{itemize}
		      \item \textbf{C. Random.nextInt() — 错误}
					  \begin{itemize}[label={}, nosep]
		      
			      \item 需要先创建Random对象
			      \item 格式：new Random().nextInt()
		      

					  \end{itemize}
		      \item \textbf{D. Math.random() — 正确}
					  \begin{itemize}[label={}, nosep]
		      
			      \item static方法，直接调用
			      \item 返回[0.0, 1.0)的double值
		      

		      
					  \end{itemize}
					  \end{itemize}
\textbf{【应背诵知识点】}

		      random()方法：
		      \begin{itemize}[label={}, nosep]
			      \item Math.random()：返回[0.0, 1.0)
			      \item Random类：nextInt(n)返回[0, n)
			      \item 生成指定范围随机数：(int)(Math.random() * n)
		      \end{itemize}
	      \end{bluebox}

	      \vspace{1em}

	\item \textbf{以下关于Java中Calendar类的说法错误的是}
	      \begin{enumerate}[label=\Alph*.]
		      \item 用于操作日期和时间
		      \item 是一个抽象类
		      \item 可以直接创建对象
		      \item 可以通过getInstance()方法获取实例
	      \end{enumerate}

	      \begin{bluebox}{答案与深度解析}
		      \textbf{答案：C. 可以直接创建对象}

		      \textbf{【核心认知框架】}

		      \textbf{【选项原理深度拆解】}
					  \begin{itemize}[label={}, leftmargin=*, nosep]


		      \item \textbf{A. 操作日期时间 — 正确}
					  \begin{itemize}[label={}, nosep]
		      
			      \item Calendar提供日期时间操作API
		      

					  \end{itemize}
		      \item \textbf{B. 抽象类 — 正确}
					  \begin{itemize}[label={}, nosep]
		      
			      \item Calendar是抽象类
		      

					  \end{itemize}
		      \item \textbf{C. 可以直接创建对象 — 错误}
					  \begin{itemize}[label={}, nosep]
		      
			      \item 抽象类不能直接new
			      \item 需要通过getInstance()
		      

					  \end{itemize}
		      \item \textbf{D. getInstance()获取实例 — 正确}
					  \begin{itemize}[label={}, nosep]
		      
			      \item Calendar cal = Calendar.getInstance();
		      

		      
					  \end{itemize}
					  \end{itemize}
\textbf{【应背诵知识点】}

		      Calendar类：
		      \begin{itemize}[label={}, nosep]
			      \item 抽象类，需通过getInstance()获取实例
			      \item JDK8+推荐使用LocalDateTime
			      \item 常用字段：YEAR、MONTH、DATE、HOUR等
			      \item 注意：MONTH从0开始（0=一月）
		      \end{itemize}
	      \end{bluebox}

	      \vspace{1em}

	\item \textbf{以下哪个方法用于格式化日期？}
	      \begin{enumerate}[label=\Alph*.]
		      \item format()
		      \item parse()
		      \item DateFormat.format
		      \item SimpleDateFormat.format()
	      \end{enumerate}

	      \begin{bluebox}{答案与深度解析}
		      \textbf{答案：D. SimpleDateFormat.format()}

		      \textbf{【核心认知框架】}

		      \textbf{【选项原理深度拆解】}
					  \begin{itemize}[label={}, leftmargin=*, nosep]


		      \item \textbf{A/B. format()/parse() — 错误}
					  \begin{itemize}[label={}, nosep]
		      
			      \item 需要通过具体类调用
		      

					  \end{itemize}
		      \item \textbf{C. DateFormat.format — 错误}
					  \begin{itemize}[label={}, nosep]
		      
			      \item DateFormat是抽象类，不能直接使用
		      

					  \end{itemize}
		      \item \textbf{D. SimpleDateFormat.format() — 正确}
					  \begin{itemize}[label={}, nosep]
		      
			      \item SimpleDateFormat是DateFormat的实现类
			      \item 格式：new SimpleDateFormat("yyyy-MM-dd").format(date)
		      

		      
					  \end{itemize}
					  \end{itemize}
\textbf{【应背诵知识点】}

		      日期格式化：
		      \begin{itemize}[label={}, nosep]
			      \item SimpleDateFormat：线程不安全
			      \item JDK8+：DateTimeFormatter（线程安全）
			      \item 常用格式：yyyy年MM月dd日 HH:mm:ss
		      \end{itemize}
	      \end{bluebox}

	      \vspace{1em}

	\item \textbf{以下关于Java中Optional类的说法错误的是}
	      \begin{enumerate}[label=\Alph*.]
		      \item 用于避免空指针异常
		      \item 可以包含值或为空
		      \item 不支持值的获取和判断
		      \item 提供了一些有用的方法
	      \end{enumerate}

	      \begin{bluebox}{答案与深度解析}
		      \textbf{答案：C. 不支持值的获取和判断}

		      \textbf{【核心认知框架】}

		      \textbf{【选项原理深度拆解】}
					  \begin{itemize}[label={}, leftmargin=*, nosep]


		      \item \textbf{A. 避免空指针 — 正确}
					  \begin{itemize}[label={}, nosep]
		      
			      \item Optional提供null安全
		      

					  \end{itemize}
		      \item \textbf{B. 可包含值或为空 — 正确}
					  \begin{itemize}[label={}, nosep]
		      
			      \item Optional.of(value) / Optional.empty()
		      

					  \end{itemize}
		      \item \textbf{C. 不支持获取和判断 — 错误}
					  \begin{itemize}[label={}, nosep]
		      
			      \item 支持！isPresent()、get()、orElse()等
		      

					  \end{itemize}
		      \item \textbf{D. 提供有用方法 — 正确}
					  \begin{itemize}[label={}, nosep]
		      
			      \item map()、filter()、orElseGet()等
		      

		      
					  \end{itemize}
					  \end{itemize}
\textbf{【应背诵知识点】}

		      Optional常用方法：
		      \begin{itemize}[label={}, nosep]
			      \item isPresent()：判断是否有值
			      \item get()：获取值（可能有NoSuchElementException）
			      \item orElse(default)：有值返回，无值返回默认值
			      \item orElseGet(supplier)：无值时执行supplier
			      \item map()：转换值
			      \item filter()：过滤值
		      \end{itemize}
	      \end{bluebox}

\end{enumerate}

\end{document}
